\chapter{Gastric Bypass}
\index{Gastric Bypass}

\section*{Definition and Overview}
Gastric bypass, formally Roux-en-Y gastric bypass, is a weight-loss (bariatric) operation that makes the stomach smaller and bypasses part of the small intestine. It is performed for people with severe obesity, usually a body mass index (BMI) of 40 or above, or 35 with serious obesity-related health conditions, who have not achieved lasting weight loss with other methods. The operation produces substantial weight loss and major improvements in diabetes, blood pressure, sleep apnoea, and quality of life. It is major surgery with real risks, and it requires lifelong changes in eating, nutrition, and follow-up. This chapter explains how gastric bypass works, its benefits, and what it involves.

\section*{Epidemiology}
Obesity affects around two-thirds of adults, and severe obesity affects a substantial minority. Bariatric surgery, including gastric bypass, is performed on thousands of people each year, and demand has risen with the obesity epidemic. Gastric bypass produces the greatest and most durable weight loss of the common bariatric procedures, averaging 60--80% of excess weight, and it leads to remission of type 2 diabetes in around 80% of people in the short to medium term. Serious complications are uncommon but real, occurring in a low single-digit percentage of operations, and long-term follow-up is essential.

\section*{Aetiology and Pathophysiology}
\begin{itemize}
    \item \textbf{The operation:} the stomach is divided into a small upper pouch, and the small intestine is connected (anastomosed) to the pouch, so food bypasses the lower stomach and the first part of the small intestine (duodenum)
    \item \textbf{Mechanisms of weight loss:} a smaller stomach restricts intake; bypassing the duodenum changes hormones, including reducing ghrelin (hunger hormone) and increasing GLP-1 (satiety hormone)
    \item \textbf{Metabolic effects:} the hormonal changes improve blood sugar control rapidly, often before significant weight loss, and treat type 2 diabetes
    \item \textbf{Malabsorption:} some malabsorption of nutrients occurs, which is why lifelong vitamin and mineral supplementation is required
    \item \textbf{Dumping syndrome:} eating sugar-rich food causes rapid gastric emptying, producing sweating, palpitations, and diarrhoea, which discourages unhealthy eating
    \item \textbf{Anaesthesia and surgery risks:} major abdominal surgery under general anaesthesia
    \item \textbf{Lifelong changes:} smaller meals, chewing well, vitamin supplements, and regular review are essential
    \item \textbf{Why it works:} combination of restriction, hormonal change, and behaviour change
\end{itemize}

\section*{Clinical Features}
The features are those of the procedure, recovery, and long-term effects.
\begin{itemize}
    \item Pre-operative assessment: medical, nutritional, and psychological evaluation, with a trial of non-surgical weight loss
    \item The operation takes 2--4 hours, usually laparoscopically, with a 1--3 night hospital stay
    \item Rapid weight loss begins within weeks and continues for 12--18 months
    \item Early satiety: eating small amounts causes fullness
    \item Diabetes remission in most people with type 2 diabetes
    \item Improved blood pressure, sleep apnoea, joint pain, and mobility
    \item The need for lifelong vitamin and mineral supplements
    \item Follow-up with the bariatric team, dietitian, and GP for life
\end{itemize}

\section*{Differential Diagnosis}
The decision for surgery involves weighing the options.
\begin{itemize}
    \item Non-surgical weight management, including diet, exercise, and medicines such as GLP-1 agonists
    \item Other bariatric procedures: gastric sleeve, gastric band, and duodenal switch
    \item The suitability of each procedure for the individual's health and preferences
    \item Whether obesity-related conditions justify the surgical risk
    \item The need to exclude medical and psychological contraindications
    \item The choice is individualised and made with the bariatric team
\end{itemize}

\section*{When to Seek Medical Care}
Anyone considering gastric bypass should have a thorough assessment by a bariatric team, including a GP, surgeon, dietitian, and psychologist. The decision is significant and should not be rushed.

\textbf{Contact your local emergency services for:}
\begin{itemize}
    \item Severe abdominal pain, fever, or vomiting in the weeks after surgery, which may indicate a leak or infection
    \item Chest pain or breathlessness, which may indicate a clot or heart problem
    \item Wound redness, discharge, or opening
    \item Inability to keep fluids down
    \item Any emergency symptoms
\end{itemize}

\section*{Diagnosis and Investigations}
\begin{itemize}
    \item \textbf{BMI and health assessment:} weight, and the presence of obesity-related conditions
    \item \textbf{Medical work-up:} blood tests, cardiac and respiratory assessment, and sleep studies where indicated
    \item \textbf{Nutritional assessment:} baseline vitamin levels, including B12, iron, and vitamin D
    \item \textbf{Psychological assessment:} eating behaviours, readiness for change, and mental health
    \item \textbf{Endoscopy:} examination of the stomach before surgery
    \item \textbf{Imaging after surgery:} X-ray contrast studies to check the join when a leak is suspected
\end{itemize}

\section*{Management}
\begin{itemize}
    \item \textbf{Before surgery:} weight loss, stopping smoking, and optimising diabetes and other conditions
    \item \textbf{The operation:} laparoscopic Roux-en-Y gastric bypass, with small incisions and a short stay
    \item \textbf{After surgery:} a staged diet from fluids to soft foods to small regular meals
    \item \textbf{Vitamins and minerals:} lifelong supplements, including B12, iron, calcium, and a multivitamin
    \item \textbf{Exercise:} regular activity to preserve muscle and support weight loss
    \item \textbf{Dietary habits:} small, frequent, protein-first meals; chew thoroughly; avoid sugary and high-fat foods
    \item \textbf{Long-term follow-up:} regular review by the bariatric team and GP, with blood tests for nutrition
    \item \textbf{Complication management:} treatment of leaks, strictures, ulcers, and nutritional deficiencies if they occur
    \item \textbf{Psychological support:} ongoing support for eating behaviour and body image
\end{itemize}

\section*{Complications}
\begin{itemize}
    \item Early surgical complications: bleeding, infection, leaks from the join, and blood clots
    \item Stricture of the join, causing difficulty eating
    \item Internal hernias and bowel obstruction
    \item Ulcers at the join
    \item Nutritional deficiencies, including B12, iron, calcium, and protein, without supplements
    \item Dumping syndrome from sugary foods
    \item Weight regain in a proportion of people
    \item Risks of anaesthesia
\end{itemize}

\section*{Prognosis and Outcomes}
Gastric bypass produces the most substantial and durable weight loss of the bariatric procedures, with most people losing 60--80% of excess weight and keeping much of it off for a decade or more. Type 2 diabetes remits in the majority, and improvements in blood pressure, sleep apnoea, mobility, and quality of life are well documented. Serious complications are uncommon, and mortality is low, but the operation demands lifelong commitment to nutrition, supplements, and follow-up. For appropriately selected people, the long-term health benefits far outweigh the risks.

\section*{Prevention and Self-Care}
\begin{itemize}
    \item Understand that surgery is a tool, not a cure, and commit to permanent lifestyle change
    \item Follow the staged diet after surgery exactly
    \item Take all vitamin and mineral supplements for life
    \item Eat small, protein-rich meals and chew slowly
    \item Avoid sugary drinks and foods that cause dumping
    \item Attend all follow-up appointments and blood tests
    \item Exercise regularly
    \item Seek prompt medical care for abdominal pain, fever, or vomiting
\end{itemize}

\section*{Sources and Further Reading}
The following reputable sources provide further information for healthcare students and patients seeking reliable, current advice about gastric bypass.
\begin{itemize}
    \item Obesity Surgery Society of Australia and New Zealand. \textit{Bariatric surgery information.} https://ossanz.com.au/
    \item Healthdirect Australia. \textit{Gastric bypass.} https://www.healthdirect.gov.au/
    \item The Royal Australasian College of Surgeons. \textit{Bariatric surgery.} https://www.surgeons.org/
    \item NHS. \textit{Gastric bypass.} https://www.nhs.uk/conditions/weight-loss-surgery/
    \item Mayo Clinic. \textit{Gastric bypass.} https://www.mayoclinic.org/
    \item World Health Organization. \textit{Obesity.} https://www.who.int/
\end{itemize}
